\section{Dante and the Hermetic Method}

James led a public discussion of the inner meaning of Dante's Divine Comedy.

Here are links to the text to be discussed:

One is the Longfellow translation chosen because it's easy to find online:

Divine Comedy Volume 1 Canto 9\footnote{\url{https://en.wikisource.org/wiki/Divine\_Comedy\_(Longfellow\_1867)/Volume\_1/Canto\_9}}

The other is the Carlyle translation which is a literal, prose translation of the text which might help in at least attempting to discern the Italian for the non-Italian speakers including myself. If anyone has a working knowledge of Dante's Italian, it would be helpful insofar as being able to see the colours of an icon will help its reader even if it doesn't automatically impart the gift of depth.

Carlyle translation\footnote{\url{https://archive.org/details/dantesdivinecome00dant/page/n29}}

James' description of the meeting follows:

I would like to take the ``summer break'' as an opportunity for some practice of the process of Mysticism-Gnosis-Magic-Hermetic Philosophy on one particular text: The Divine Comedy.

I would like to propose that, for anyone interested, we gather to meditate as a group on Canto IX of the Inferno: the native canto of Dante's plea that: ``Ye that are of good understanding, note the teaching that is hidden under the veil of the strange lines.''

I would further like to propose that the format of this gathering is not one of scholastic discussion on the ``relevance'' of the poem, or introduction to the poem, or even as a practice of ``intercontextualization'' with other works we have read, but as a participation in the ``method'' of those other works in cooperation with Dante's lines. In other words, to go from knowledge of Dante and his milieu of tradition, to gnosis of Dante through direct mystical experience with the poem.

To that end, I would like to propose a period of calm reading followed by quiet meditation and then cooperative discussion as a starting template. I am no expert in confraternities attempting to achieve depth so I am open to suggestions on a salient format.

As usual, there is no ``lesson'' that I will be attempting to ``teach'' --- I proclaim no masterhood over the poem. But, rather, I wish to participate as a fellow pilgrim hoping that what little I have experienced may be combined with those of others so that we may, perhaps, contribute to the Hermetic Current through our interpenetration with the sublime poet.

Thus, the objective would be to, as always, somehow increase or practice our heightened awareness and hopefully bring us closer to that state of purgation, illumination, and union which the Poet-Pilgrim arduously accomplished in his Comedy. If we are able to achieve gnosis of the text, we may even witness a state of magic if it somehow facilitates a miracle to be achieved in our lives. Hopefully, if magic does confirm the gnosis we have received through this mystical experience of the lines, then perhaps it will be possible for a cooperation on a ``book'' to be achieved---a codification and summary of these previous steps.

In a sense, I propose an imitation of Dante's method which is the path of the lozenge. He experienced the reality of the Religious Sense in pure mystical experience that had descended onto his horizontal plane and he undertook this task of mirroring that triangle into a lozenge by taking the totality of his experience and concentrating it effortlessly (poetry) into a singularity-mirror of that divine singularity above (similar to the transition of Teilhard de Chardin's Alpha to Omega). I wish to``repeat'' that process by taking the singularity of Dante's poem,expanding our gnosis of it and once again synthesizing into a ``book''so as to form two diamonds: one greater, one lesser.

Perhaps under a microscope, we can look at Dante's own mystical experience as the interaction with the Tree of Life which he then reiterates in his poetry through the``particular examples'' of the four worlds forming layer after layer of poetic polyvalence similar to how the Tree can be seen across the four worlds or how the suits of the Tarot continue the world-as-magic task of the Major Arcana. Thus, in the same way that the final Sefirah can be the starting point of the next Tree, so can Dante's work be taken as our point to practice, in microcosm, what he so perfectly accomplished.

To that end, analogous to how imitation of Mary is also imitation of Christ, to undergo the task of understanding through gnosis the Divine Comedy which so perfectly bears the Christian event into the world of letters is also a work of understanding the DIVINE COMEDY---the salvation history of man as the prodigal son---which the Divine Comedy mirrors.

Thank you, and, simultaneously, I hope that this will contribute to the particular interests of each friend involved.


\flrightit{Posted on 2019-06-12 by Admin}

